\documentclass[11pt]{article}
\usepackage[utf8]{inputenc}
\usepackage[T1]{fontenc}
\usepackage[french]{babel}
\usepackage{lmodern}
\usepackage[a4paper,margin=2.2cm]{geometry}
\usepackage{amsmath}
\usepackage{setspace}
\usepackage{enumitem}

\setlength{\parindent}{0pt}
\setlength{\parskip}{0.6em}
\setlist[itemize]{leftmargin=1.2em}

\begin{document}


\begin{center}
{\LARGE Activité : Conception d’un plan d’échantillonnage\\[0.4em]}
{\Large \textit{Jelly Blubbers}}
\end{center}

\textbf{Échantillon de convenance —}
Sélectionnez 5 Jelly Blubbers faciles à collecter – ceux qui sont les plus simples à repérer sur la feuille.
Notez leurs longueurs et calculez la moyenne.\\[-0.4em]

\rule{1.6cm}{0.4pt}\hspace{0.6em}
\rule{1.6cm}{0.4pt}\hspace{0.6em}
\rule{1.6cm}{0.4pt}\hspace{0.6em}
\rule{1.6cm}{0.4pt}\hspace{0.6em}
\rule{1.6cm}{0.4pt}\hfill
$\bar{x} = $ \rule{2.2cm}{0.4pt}

\textbf{Échantillon raisonné (de jugement) —}
Sélectionnez 5 Jelly Blubbers qui, selon votre jugement, sont représentatifs de la population sur cette feuille.
Notez leurs longueurs et calculez la moyenne.\\[-0.4em]

\rule{1.6cm}{0.4pt}\hspace{0.6em}
\rule{1.6cm}{0.4pt}\hspace{0.6em}
\rule{1.6cm}{0.4pt}\hspace{0.6em}
\rule{1.6cm}{0.4pt}\hspace{0.6em}
\rule{1.6cm}{0.4pt}\hfill
$\bar{x} = $ \rule{2.2cm}{0.4pt}

\textbf{Échantillon aléatoire simple —}
Générez cinq entiers aléatoires entre 1 et 100. Repérez les Jelly Blubbers correspondants et enregistrez leurs longueurs.
Calculez la moyenne.\\[-0.4em]

\rule{1.6cm}{0.4pt}\hspace{0.6em}
\rule{1.6cm}{0.4pt}\hspace{0.6em}
\rule{1.6cm}{0.4pt}\hspace{0.6em}
\rule{1.6cm}{0.4pt}\hspace{0.6em}
\rule{1.6cm}{0.4pt}\hfill
$\bar{x} = $ \rule{2.2cm}{0.4pt}

\textbf{Échantillonnage stratifié —}
Supposez que l’on secoue les Jelly Blubbers, ce qui fait couler les plus gros et remonter les plus petits.
On obtient alors une feuille stratifiée où les tailles sont divisées en 5 strates.
Choisissez un Blubber au hasard dans chaque strate (en générant un entier aléatoire),
repérez-les, notez leurs longueurs et calculez la moyenne.\\[-0.4em]

\rule{1.6cm}{0.4pt}\hspace{0.6em}
\rule{1.6cm}{0.4pt}\hspace{0.6em}
\rule{1.6cm}{0.4pt}\hspace{0.6em}
\rule{1.6cm}{0.4pt}\hspace{0.6em}
\rule{1.6cm}{0.4pt}\hfill
$\bar{x} = $ \rule{2.2cm}{0.4pt}

\textbf{Échantillonnage systématique —}
Utilisez un système pour sélectionner l’échantillon en utilisant la page originale des Blubbers pour sélectionnez cinq Blubbers.
Générez un entier aléatoire entre 1 et 20 : c’est le premier Blubber.
Ajoutez ensuite 20 au nombre précédent pour trouver les quatre autres.
Notez leurs longueurs et calculez la moyenne.\\[-0.4em]

\rule{1.6cm}{0.4pt}\hspace{0.6em}
\rule{1.6cm}{0.4pt}\hspace{0.6em}
\rule{1.6cm}{0.4pt}\hspace{0.6em}
\rule{1.6cm}{0.4pt}\hspace{0.6em}
\rule{1.6cm}{0.4pt}\hfill
$\bar{x} = $ \rule{2.2cm}{0.4pt}

\textbf{Échantillonnage par grappes —}
Les échantillons par grappes sélectionnent des groupes entiers (généralement basés sur la localisation).
Générez un entier aléatoire $r$ entre 1 et 20 puis multipliez-le par 5.
Votre échantillon sera ce nombre ($5r$) et les quatre nombres qui le précèdent.
(Ex. : $r = 12$, $5r = 60$ ; l’échantillon est 60, 59, 58, 57, 56.)
Notez leurs longueurs et calculez la moyenne.\\[-0.4em]

\rule{1.6cm}{0.4pt}\hspace{0.6em}
\rule{1.6cm}{0.4pt}\hspace{0.6em}
\rule{1.6cm}{0.4pt}\hspace{0.6em}
\rule{1.6cm}{0.4pt}\hspace{0.6em}
\rule{1.6cm}{0.4pt}\hfill
$\bar{x} = $ \rule{2.2cm}{0.4pt}



\end{document}
